\newpage
\renewcommand\thesection{1}
\renewcommand\thesubsection{\arabic{subsection}}

\setcounter{subsection}{0}
\counterwithin{subsection}{section}

\pagenumbering{arabic} 

\begin{center}
    \section*{BAB I \\ PENDAHULUAN}
    \addcontentsline{toc}{section}{BAB I \quad PENDAHULUAN}
\end{center}

\subsection{Latar Belakang Masalah}

    \hspace{0.5cm}Korupsi masih menjadi salah satu permasalahan serius di Indonesia karena tidak hanya menimbulkan kerugian negara, tetapi juga merusak kepercayaan masyarakat, melemahkan tata kelola pemerintahan, dan mencederai nilai keadilan sosial. Praktik korupsi dapat terjadi dalam berbagai bentuk, mulai dari penyalahgunaan jabatan, gratifikasi, suap, penggelapan dana, hingga perilaku tidak jujur dalam kehidupan sehari-hari. Permasalahan ini menunjukkan bahwa korupsi bukan hanya persoalan hukum, tetapi juga persoalan moral, karakter, dan kebiasaan sosial masyarakat \citep{nainggolan2024}. Dalam lingkup yang lebih kecil, perilaku seperti mencontek, titip absen, plagiarisme, dan manipulasi akademik juga dapat dipandang sebagai bentuk awal dari perilaku tidak berintegritas yang berpotensi berkembang menjadi budaya koruptif \citep{devina2026}.

    \hspace{0.5cm}Indonesia dikenal sebagai negara yang masyarakatnya memiliki tingkat religiusitas yang tinggi. Nilai-nilai agama seharusnya menjadi dasar penting dalam membentuk perilaku jujur, bertanggung jawab, dan berintegritas. Setiap agama pada dasarnya mengajarkan umatnya untuk menjauhi tindakan mengambil hak orang lain secara tidak sah. Namun, tingginya religiusitas masyarakat belum selalu tercermin dalam perilaku antikorupsi. Fenomena ini menunjukkan adanya kesenjangan antara pengetahuan agama yang dipahami secara normatif dengan penerapannya dalam kehidupan sehari-hari. Nilai agama sering kali berhenti pada aspek ritual dan simbol keagamaan, tetapi belum sepenuhnya menjadi pedoman moral dalam menghadapi tekanan sosial, kepentingan pribadi, dan dilema etis di masyarakat \citep{siagian2025,zahron2026}.

    \hspace{0.5cm}Berdasarkan kondisi tersebut, peneliti tertarik untuk melakukan kegiatan studi lapangan mengenai pandangan agama dalam penanaman moral antikorupsi kepada masyarakat. Kegiatan ini dilakukan melalui wawancara dengan tokoh agama dan penyebaran kuesioner kepada masyarakat untuk mengetahui sejauh mana nilai keagamaan dipahami serta diterapkan dalam membangun sikap antikorupsi. Melalui kegiatan ini, peneliti ingin melihat bagaimana ajaran agama memandang tindakan korupsi, faktor apa saja yang menyebabkan praktik korupsi tetap terjadi meskipun masyarakat memiliki pemahaman agama, serta bagaimana peran tokoh agama, lembaga keagamaan, dan lingkungan sosial dalam menanamkan nilai kejujuran, tanggung jawab, dan integritas.

\subsection{Rumusan Masalah}
 
\begin{enumerate}
    \item Bagaimana pandangan ajaran agama terhadap tindakan korupsi dan nilai integritas dalam kehidupan bermasyarakat?
    \item Apa faktor-faktor yang menyebabkan praktik korupsi tetap terjadi meskipun masyarakat Indonesia memiliki tingkat religiusitas yang tinggi?
    \item Bagaimana peran pendidikan agama, tokoh agama, dan lingkungan sosial dalam menanamkan moral antikorupsi kepada masyarakat, khususnya generasi muda?
\end{enumerate}

\subsection{Tujuan dan Manfaat}
\begin{minipage}{1\linewidth}
    \subsubsection{Tujuan}
     Berdasarkan rumusan masalah di atas, tujuan kegiatan ini adalah untuk memahami pandangan berbagai agama terhadap tindakan korupsi serta nilai kejujuran dan integritas dalam kehidupan sosial. Selain itu, kegiatan ini juga bertujuan untuk mengidentifikasi faktor-faktor yang menyebabkan praktik korupsi tetap terjadi di tengah masyarakat religius serta mengetahui peran tokoh agama dan lingkungan sosial dalam menanamkan moral antikorupsi kepada masyarakat.

    \subsubsection{Manfaat}
    Melalui kegiatan ini, manfaat yang dapat diambil yaitu :
    \subsubsection{Secara teoritis}
    Kegiatan ini dapat membantu peneliti memahami hubungan antara nilai keagamaan, moralitas, dan perilaku antikorupsi dalam masyarakat. Hasil kegiatan ini juga dapat menjadi bahan kajian mengenai bagaimana ajaran agama berperan dalam membentuk karakter, integritas, dan kesadaran moral seseorang dalam menghadapi tindakan koruptif.
    \subsubsection{Secara praktis}
    Secara praktis, hasil kegiatan ini dapat memberikan pemahaman kepada masyarakat mengenai pentingnya menerapkan nilai agama secara nyata dalam kehidupan sehari-hari, khususnya dalam menolak tindakan korupsi. Kegiatan ini juga dapat menjadi masukan bagi lembaga pendidikan, tokoh agama, dan lingkungan sosial agar pembinaan moral tidak hanya bersifat normatif, tetapi juga lebih kontekstual dan mampu membekali masyarakat dalam menghadapi dilema moral di kehidupan nyata.
\end{minipage}

    
 \subsection{Waktu dan Lokasi Studi Lapangan}

     \begin{minipage}{1\linewidth}
     \begin{minipage}{0.2\linewidth}
         \subsubsection{Waktu}
     \end{minipage}
     \begin{minipage}[t]{0.8\linewidth}
         \tabto{0.5cm} : 15 Maret 2026 - 11 April 2026
     \end{minipage}
 \end{minipage}

 \bigskip \medskip
 \begin{minipage}{1\linewidth}
     \begin{minipage}{0.2\linewidth}
         \subsubsection{Tempat}
     \end{minipage}
     \begin{minipage}[t]{0.8\linewidth}
         \tabto{0.5cm} : 
         \begin{minipage}[t]{1\linewidth}
             Vihara Ekayana \\
             Jl. Mangga 2 No.8, RT.8/RW.8, Duri Kepa, Kec. Kebon Jeruk, Kota Jakarta Barat, Daerah Khusus Ibukota Jakarta 11510 \\

             GPKdI Blessing Community Jelambar \\
             Jl. Jelambar Madya Utara No.145A, RT.8/RW.2, Jelambar, Kec. Grogol Petamburan, Kota Jakarta Barat, Daerah Khusus Ibukota Jakarta 11460 \\

             Komunitas Mahasiswa Katolik Jakarta Barat \\
             Jl. Gelong Baru Utara II No.36C 15, RT.15/RW.7, Tomang, Kec. Grogol Petamburan, Kota Jakarta Barat, Daerah Khusus Ibukota Jakarta 11440 \\

             Pasraman Chandra Prabha \\
             Jl. Indraloka Raya No.1, RT.7/RW.10, Jelambar, Kec. Grogol Petamburan, Kota Jakarta Barat, Daerah Khusus Ibukota Jakarta 11460 \\

             Penyuluh Agama Islam Kec. Grogol Petamburan \\
             Jl. Anyar Raya No.1, RT.1/RW.10, Wijaya Kusuma, Kec. Grogol Petamburan, Kota Jakarta Barat, Daerah Khusus Ibukota Jakarta 11460 \\
         \end{minipage}
     \end{minipage}
 \end{minipage}


